
% ============================================================
% Supplementary Training-Budget Analysis
% ============================================================

\begin{revisedblock}
This appendix provides supplementary results on the
sensitivity of DQN and A2C reliability to training
budget. Intermediate model checkpoints corresponding
to nominal budgets of 50,000, 100,000, and 200,000
interactions were evaluated alongside the final
400,000-interaction checkpoints.

The evaluation used site1, site3, and site5, five
independent training seeds, four supplier-delay
scenarios, and ten evaluation episodes per seed and
scenario. The intermediate-checkpoint analysis
comprised 90 model evaluations and 3,600 held-out
episodes. Each checkpoint was evaluated using its
corresponding saved observation-normalisation
statistics.

Table~\ref{tab:appendix_rl_budget} reports the
results for site5, where reliability differences
were most evident. Mean EENS is calculated across
five training seeds and four delivery scenarios.

Differences are measured relative to the final
checkpoint. The exploratory 95\% confidence
intervals are based on paired training-seed results.
\end{revisedblock}


\begin{table}[htbp]
\centering
\small
\caption{\revised{Training-budget sensitivity for
DQN and A2C on site5. Mean EENS and differences
are in kWh per 360-hour episode. A positive
difference indicates higher EENS than the final
checkpoint.}}
\label{tab:appendix_rl_budget}
\begin{tabular}{llrrl}
\toprule
\textbf{Algorithm} &
\textbf{Budget} &
\textbf{Mean EENS} &
\textbf{Difference} &
\textbf{95\% CI} \\
\midrule
\revised{DQN} & \revised{50k} &
\revised{84.8580} & \revised{+3.9729} &
\revised{$[-3.2409,\ 11.1866]$} \\
\revised{DQN} & \revised{100k} &
\revised{81.0724} & \revised{+0.1872} &
\revised{$[-1.0218,\ 1.3963]$} \\
\revised{DQN} & \revised{200k} &
\revised{80.9840} & \revised{+0.0988} &
\revised{$[-0.9470,\ 1.1447]$} \\
\revised{DQN} & \revised{400k} &
\revised{80.8852} & \revised{--} &
\revised{Reference} \\
\midrule
\revised{A2C} & \revised{50k} &
\revised{83.7292} & \revised{-1.0674} &
\revised{$[-2.6841,\ 0.5492]$} \\
\revised{A2C} & \revised{100k} &
\revised{84.3715} & \revised{-0.4252} &
\revised{$[-2.3601,\ 1.5097]$} \\
\revised{A2C} & \revised{200k} &
\revised{84.0428} & \revised{-0.7539} &
\revised{$[-2.4176,\ 0.9099]$} \\
\revised{A2C} & \revised{400k} &
\revised{84.7967} & \revised{--} &
\revised{Reference} \\
\bottomrule
\end{tabular}
\end{table}

\begin{revisedblock}
For DQN, mean EENS on site5 decreases from
84.8580 kWh at 50,000 interactions to
81.0724 kWh at 100,000, followed by smaller
numerical improvements at subsequent checkpoints.
A2C does not exhibit a consistent improvement
with additional interactions.

Across all three sites, four delivery scenarios,
and three intermediate checkpoints, none of the
72 exploratory paired 95\% confidence intervals
for EENS differences relative to the final
checkpoint excludes zero. Site1 records zero
EENS for both algorithms at all checkpoints,
while site3 provides limited additional
discrimination.

These results describe training-budget sensitivity
rather than statistically established convergence.
The final 400,000-interaction checkpoints were
retained for the main algorithm comparison;
intermediate checkpoints were not selected based
on held-out evaluation performance.

Equivalent intermediate-checkpoint evaluations
were not available for the historical RLInv and A7
models. Consequently, this supplementary analysis
does not establish relative learning efficiency
across all four RL algorithms.
\end{revisedblock}
